\section{Kiến trúc kho dữ liệu}\label{sec:kien-truc}
Kiến trúc gồm sáu lớp (Hình \ref{fig:pipeline}): hai hệ thống nguồn, staging, ODS, kho dữ liệu hình sao, các data mart theo chủ đề và Power BI; lớp giám sát \code{etl\_log} ghi lại mỗi lần chạy, dòng bị loại và kết quả kiểm thử. Python điều phối pipeline, mọi biến đổi chính viết bằng SQL procedure trên PostgreSQL. Vai trò từng lớp: staging giữ bản sao thô để nạp tăng dần và đối soát; ODS làm sạch và chuẩn hóa mà không xóa dữ liệu; kho tổ chức theo mô hình chiều với khóa thay thế và SCD; data mart tổng hợp sẵn cho từng nhóm dashboard.
\hinh{3-4-1_pipeline.png}{Kiến trúc các lớp và pipeline ETL (bản nháp tự động)}{fig:pipeline}{1.0}
Số bảng và số dòng của từng lớp ở trạng thái cuối được thống kê ở Bảng \ref{tab:layers} (mục 
ef{sec:du-lieu-nguon}).
\todo{Yêu cầu của môn: sơ đồ kiến trúc \textbf{có logo công nghệ ở mỗi phần} (PostgreSQL, Python, Power BI, GitHub...), nêu các hoạt động ETL, dimension, fact và tên các dashboard. Vẽ lại bằng draw.io từ các file docs/diagrams/kien-truc-cac-lop.drawio và pipeline-etl.drawio, xuất PNG vào report/figures/.}

\subsection{Kiến trúc theo lớp: thành phần, công nghệ, số bảng và số dòng}\label{sec:kien-truc-lop}
Hình \ref{fig:pipeline} ở trên là sơ đồ luồng; để thấy \emph{cấu trúc} của hệ thống, Hình \ref{fig:kien-truc-lop} vẽ lại kiến trúc theo lớp: các khối xếp tầng từ lớp trình bày (Power BI, bảy dashboard) xuống lớp nguồn (hai hệ thống \code{oltp} và \code{crm}), mỗi lớp liệt kê thành phần, công nghệ, số bảng và số dòng. Hai thanh ngang ở dưới là các thành phần xuyên suốt: giám sát \code{etl\_log} (6 bảng) và điều phối bằng Python với mã nguồn quản lý trên GitHub. Nhóm vẽ theo lớp thay vì theo luồng vì kiến trúc trả lời câu hỏi \emph{hệ thống gồm những gì và mỗi phần dùng công nghệ nào}, còn câu hỏi \emph{dữ liệu đi từ đâu đến đâu} đã có sơ đồ data flow ở mục \ref{sec:dataflow} và sơ đồ các job ở mục \ref{sec:etl-jobs}.
\hinh{3-2_kien_truc_cac_lop.png}{Kiến trúc kho dữ liệu theo lớp (bản nháp; ô công nghệ là ô chữ, nhóm thay bằng logo chính thức)}{fig:kien-truc-lop}{1.0}
Bảng \ref{tab:p8-layers} bổ sung số bảng, số dòng ở trạng thái cuối và tần suất cập nhật của từng lớp. Công nghệ chính là PostgreSQL 18 (lưu trữ và biến đổi bằng procedure PL/pgSQL), Python 3.11 (điều phối, nạp nguồn, xuất CSV), Power BI Desktop (mô hình ngữ nghĩa và dashboard) và Git/GitHub (quản lý mã nguồn, mỗi giai đoạn một nhánh).
\begin{table}[htbp]\centering\scriptsize\caption{Các lớp của kiến trúc: thành phần, công nghệ, số bảng, số dòng và tần suất (trạng thái cuối)}\label{tab:p8-layers}
\begin{tabular}{p{1.5cm}p{3.0cm}p{3.0cm}p{1.4cm}p{1.6cm}p{2.6cm}p{2.6cm}}\toprule
\textbf{Lớp} & \textbf{Thành phần} & \textbf{Công nghệ} & \textbf{Số bảng} & \textbf{Số dòng} & \textbf{Tần suất} & \textbf{Ghi chú}\\\midrule
Nguồn & oltp (9 bảng), crm (2 bảng) & PostgreSQL 18; nạp từ 11 file CSV bằng Python (COPY) & 11 & 1.559.764 & mô phỏng cập nhật theo tháng (ban đầu đến 31/12/2017, rồi 11 đợt) & hai hệ thống độc lập, chỉ gặp nhau qua seller\_id\\
Staging & stg\_* (bản sao thô) & PostgreSQL; extract bằng Python + SQL & 11 & 1.559.764 & mỗi lần chạy pipeline (full hoặc incremental) & thêm mã lô nạp và thời điểm nạp để truy vết\\
ODS & ods.* (dữ liệu sạch, cờ chất lượng) & PostgreSQL procedure (PL/pgSQL) & 18 & 678.183 & mỗi lần chạy pipeline & dòng lỗi ghi \code{rejected\_rows}, không xoá âm thầm\\
Kho dữ liệu & dw: 10 dimension + 6 fact & PostgreSQL procedure; khoá thay thế, SCD1/SCD2 & 16 & 614.223 & mỗi lần chạy pipeline; nạp tăng dần theo watermark & 11 view đóng vai cho \code{dim\_date}\\
Data mart & mart\_sales, mart\_fulfillment, mart\_cx, mart\_acquisition & PostgreSQL view và materialized view & 17 view (4 materialized) & -- & làm mới materialized view mỗi lần chạy & tổng hợp sẵn cho dashboard D1--D6\\
Giám sát & etl\_log: job\_run, rejected\_rows, dq\_check... & PostgreSQL; Python ghi log & 6 & 2.637 & ghi ở mọi bước & nguồn cho dashboard D7\\
BI & 37 file CSV, mô hình ngữ nghĩa, 7 dashboard & Python xuất CSV; Power BI Desktop (Import, DAX) & -- & -- & làm mới sau mỗi lần xuất CSV & Publish to web\\
\bottomrule\end{tabular}\end{table}

\todo{Thay các ô chữ công nghệ trong Hình \ref{fig:kien-truc-lop} bằng logo chính thức (PostgreSQL, Python, Power BI, GitHub) khi vẽ lại bằng draw.io từ docs/diagrams/kien-truc-cac-lop.drawio; giữ nguyên tên file hình 3-2\_kien\_truc\_cac\_lop.png.}
